\chapter{Introduction}
\section{Research Background}
\subsection{Current Development of Global and Domestic Option Markets}

Since the Chicago Board Options Exchange (CBOE) officially listed standardized options for trading in 1973, the global financial derivatives market has undergone a profound evolution from a single-product market to a multi-layered, highly complex system. In mature international markets, options are no longer merely basic risk-hedging tools, but have evolved into core instruments for asset price discovery, tail risk management, and strategic gaming. In recent years, with the intensification of global macroeconomic uncertainty and frequent geopolitical conflicts, the demand among institutional investors to guard against extreme tail risks has become increasingly urgent, pushing global option trading volumes to continuous record highs. In particular, the massive popularity of ultra-short-term products represented by Zero-Days-To-Expiration (0DTE) options in the U.S. stock market has not only profoundly altered the microstructure of market liquidity but also posed more stringent theoretical challenges to the instantaneous response capabilities of pricing models at extremely short time scales, as well as their ability to accurately capture extreme volatility skew. Consequently, the limitations of traditional constant volatility models and linear stochastic volatility models have become increasingly prominent.

At the same time, the Chinese options market, as an important emerging force in the global derivatives landscape, is in a critical window of leapfrog development. In terms of exchange-traded markets, since the successful pilot of SSE 50 ETF options in 2015, China has rapidly built an equity option matrix covering core indices such as the CSI 300, CSI 500, and CSI 1000, as well as characteristic indices like the STAR 50 and ChiNext Index. The commodity options sector has achieved full-category coverage, ranging from non-ferrous metals and the ferrous industry chain to agricultural products and new energy raw materials. This provides a rich set of risk-hedging tools for real-economy enterprises and serves as an important financial pillar for their robust operations. In the over-the-counter (OTC) derivatives sector, the OTC options market—with the China Securities inter-institution quotation system as its core link—has developed steadily. The increasing abundance of personalized and customized risk solutions has further perfected the risk management functions of the multi-layered capital market.

\subsection{Evolution of Option Pricing Models and Computational Methods}

The evolution of option pricing theory constitutes the core logical thread in the development of modern financial engineering. Since Black, Scholes, and Merton proposed the classic BSM pricing model, the global financial derivatives market has undergone a profound evolution from a single-product market to a multi-layered, highly complex system. In mature international markets, options are no longer merely basic risk-hedging tools, but have evolved into core instruments for asset price discovery, tail risk management, and strategic gaming. In recent years, with the intensification of global macroeconomic uncertainty and frequent geopolitical conflicts, the demand among institutional investors to guard against extreme tail risks has become increasingly urgent, pushing global option trading volumes to continuous record highs. In particular, the massive popularity of ultra-short-term products represented by Zero-Days-To-Expiration (0DTE) options in the U.S. stock market has not only profoundly altered the microstructure of market liquidity but also posed more stringent theoretical challenges to the instantaneous response capabilities of pricing models at extremely short time scales, as well as their ability to accurately capture extreme volatility skew. Consequently, the limitations of traditional constant volatility models and linear stochastic volatility models have become increasingly prominent.

In the early explorations of stochastic volatility models, the Heston model (also known as the 1/2 model) became the standard benchmark for global derivatives pricing due to its elegant characterization of volatility's mean-reverting property and the closed-form expression of the underlying log-price characteristic function. Subsequently, to address the complex volatility dynamics in interest rate and foreign exchange markets, the SABR model introduced a stochastic volatility term and a correlation structure between volatility and the underlying price, further strengthening its ability to capture the shape of the implied volatility curve, thereby becoming the mainstream model for OTC derivatives pricing. However, when dealing with the steep implied volatility curves of ultra-short-term options and the explosive jumps in volatility under extreme market conditions, the linear diffusion structure of the traditional 1/2 model reveals rigid theoretical limitations—its linear drift term struggles to explain the nonlinear, accelerated mean-reverting characteristics of volatility during extreme crises, nor can it adequately capture the heavy-tailed distribution of volatility in highly volatile markets.

Against this backdrop, the 3/2 stochastic volatility model has gradually entered the core research focus of both academia and industry. Unlike the Heston model, the 3/2 model assumes that the instantaneous variance process follows a nonlinear power-law diffusion process, where the diffusion term is proportional to the instantaneous variance raised to the power of 3/2, and the drift term exhibits a nonlinear mean-reverting structure. This nonlinear specification endows the model with tremendous flexibility, enabling it to more sensitively reflect volatility clustering effects under extreme market stress and more accurately fit the volatility smile across different maturities and strike prices. Consequently, it demonstrates fitting precision significantly superior to traditional affine models in the pricing practices of cross-asset classes such as equities, foreign exchange, and cryptocurrencies.

Accompanied by the increasing dimensionality and structural complexity of pricing models, the evolution of numerical computational methods has simultaneously driven the expansion of financial engineering practices. Early Finite Difference Methods (FDM) provided robust pricing pathways for path-dependent derivatives like American options by discretizing partial differential equations (PDEs); however, under multi-factor stochastic volatility models, they often face the challenge of the ``curse of dimensionality,'' with computational complexity growing exponentially as model dimensions increase. Monte Carlo Simulation, with its powerful versatility, has become the mainstream tool for handling high-dimensional models and path-dependent exotic options, but it suffers from slow convergence rates, and computational costs rise significantly as precision requirements increase. To balance pricing efficiency and precision, Conditional Monte Carlo introduces variance reduction techniques to decouple the underlying price process from the volatility process, greatly optimizing the simulation effectiveness and convergence speed of stochastic volatility models. In recent years, Fourier Transform Methods based on characteristic functions, combined with numerical algorithms like the Fast Fourier Transform (FFT) and Fourier-Cosine (COS) methods, have achieved high-speed analytical solutions for complex log-prices. This has laid a solid technical foundation for the real-time pricing and full-surface calibration of non-affine stochastic volatility models such as the 3/2 model.

\subsection{Research Questions}

Although the 3/2 model possesses significant theoretical advantages in characterizing extreme market volatility and fitting volatility surfaces, its non-affine mathematical structure also brings severe numerical implementation challenges to option pricing. Existing research has proposed various option pricing methods for the 3/2 model, encompassing three major paradigms: time-stepping Monte Carlo simulations, exact/near-exact terminal simulations, and Fourier time-domain pricing. These cumulatively form nine mainstream implementation schemes; however, three core unresolved issues remain in current research and practice:

First, existing performance evaluations of various pricing methods under the 3/2 model are mostly limited to idealized academic benchmark cases in stable market environments, lacking systematic testing of the methods' full-scenario numerical stability. The majority of nominally ``exact'' pricing methods can only achieve theoretical validation under the benchmark parameters given in literature. Under extreme parameter combinations—such as high volatility-of-volatility (vol-of-vol), strong price-volatility correlation, and extremely short time-to-maturity—problems like numerical overflow of special functions, distribution fitting failures, and non-numerical calculation results (NaN/Inf) are prevalent. The engineering applicability boundaries of these methods have yet to be clearly defined.

Second, there is a lack of empirical re-validation for pricing methods under genuinely extreme market conditions; theoretical precision in idealized examples does not equate to practical performance in real markets. The actual parameters of extremely volatile markets, such as cryptocurrencies, often exceed the parameter ranges of academic benchmark cases. Structural mutations of model parameters during extreme market conditions further amplify the performance divergence among different pricing methods. Existing research has not yet clarified the pricing accuracy, computational efficiency, and robustness of different methods under real market extreme shocks, failing to provide a reliable basis for method selection in derivatives pricing and risk management during extreme conditions.

Third, empirical adaptation research of the 3/2 model in high-volatility markets remains insufficient, lacking a dynamic evolution analysis of model parameters throughout the full cycle of extreme events. Current research focuses on static parameter calibration and the derivation of pricing formulas for the 3/2 model. It has neither constructed a comprehensive analytical framework of ``long-term fixed-parameter estimation - time-varying parameter cross-sectional calibration'' for extreme crash events, nor validated the model's ability to capture the evolution of the volatility smile under extreme conditions within a pure diffusion framework. The practical value of the model in highly volatile markets has yet to receive adequate empirical support.

Based on this, the core research questions of this paper focus on three levels: Firstly, to systematically clarify the theoretical logic, error sources, and convergence characteristics of the nine mainstream option pricing methods under the 3/2 model, and to quantitatively evaluate their pricing accuracy, computational efficiency, and full-scenario numerical stability through multi-scenario numerical experiments, thereby forming a clear hierarchical method performance system. Secondly, to execute an extreme-scenario re-evaluation of pricing methods based on model parameters calibrated from real extreme market events, comparing the performance differences between idealized cases and real markets to define the engineering applicability boundaries of different methods. Thirdly, to validate the empirical adaptability of the 3/2 model in extreme cryptocurrency market conditions, identifying the time-varying patterns of core model parameters throughout the full cycle of extreme events, thus providing theoretical support and empirical foundations for derivatives pricing and risk management in high-volatility environments.

\subsection{Research Objectives and Main Contents}

Addressing the core research questions mentioned above, this paper takes European option pricing under the 3/2 stochastic volatility model as its central research object. Through systematic research encompassing theoretical construction, numerical experiments, and empirical analysis, this study comprehensively evaluates the full-scenario performance of different pricing methods, clarifies their applicability boundaries, and simultaneously validates the empirical adaptability of the 3/2 model in extreme market environments. The core research objectives and main contents of this paper are divided into the following three aspects:

First, to construct a comprehensive methodological framework for European option pricing under the 3/2 model, accomplishing the theoretical derivation and numerical implementation of various methods. This paper categorizes the nine mainstream pricing methods into three major paradigms: time-stepping conditional Monte Carlo methods, exact and near-exact terminal simulation methods, and Fourier frequency-domain pricing methods. It systematically derives the core theoretical logic, numerical implementation processes, error characteristics, and convergence properties of each method paradigm. This addresses the problem of scattered methodological frameworks and unclear theoretical boundaries in existing research, building a unified theoretical structure and reproducible implementation schemes for subsequent numerical comparisons and empirical testing.

Second, to execute a systematic performance evaluation of different pricing methods through multi-scenario numerical experiments. This paper designs seven sets of test cases covering both conventional markets and extreme parameter environments. Using numerical stability, pricing accuracy, and computational efficiency as core evaluation metrics, it quantitatively compares the performance of the nine methods under different parameter combinations. This identifies the advantageous scenarios and core defects of each method, forming a clear hierarchical performance system. It resolves the issues of unilateral method performance evaluation and insufficient extreme-scenario robustness testing found in existing literature, providing a quantitative basis for method selection across different application scenarios.

Third, to perform empirical calibration of the 3/2 model and extreme-scenario re-evaluation of pricing methods based on real extreme market events. Taking the extreme Bitcoin market flash crash of October 2025 as the research sample, this paper constructs a two-step empirical framework of ``long-term fixed-parameter time-series estimation - time-varying parameter cross-sectional calibration.'' It identifies the time-varying patterns of the core 3/2 model parameters throughout the crash cycle and validates the model's ability to fit the evolution of the volatility smile during extreme conditions. Simultaneously, based on the real market parameters obtained from calibration, it constructs extreme scenario test cases to complete a performance re-evaluation of the pricing methods. By conducting a horizontal comparison with the academic benchmark results, it clarifies the core differences in method performance between idealized environments and real markets, ultimately forming method selection guidelines for extreme market environments, thus bridging the gap between theoretical examples and real-world markets in current research.

\subsection{Structure of the Thesis}

This thesis consists of seven chapters, organized along a logical thread of ``Theoretical Foundation - Method Construction - Numerical Experiments - Empirical Testing - Conclusions and Outlook.'' The core contents and structural arrangements of each chapter are as follows:

Chapter One is the Introduction. This chapter systematically expounds on the research background and significance of this study, reviews the developmental trajectory of global and domestic option pricing theories, defines the core research questions, objectives, and main contents concerning the 3/2 stochastic volatility model, and outlines the overall structural arrangement of the thesis.

Chapter Two is the Literature Review. This chapter systematically combs through the theoretical evolution of stochastic volatility models and the current research status of exact simulation pricing methods. It summarizes applied research of the 3/2 model in financial derivatives markets and relevant domestic research, ultimately clarifying the deficiencies in existing studies through literature review to establish the academic positioning and marginal contributions of this paper.

Chapter Three covers the Specification of the 3/2 Stochastic Volatility Model. This chapter details the core theoretical setup of the 3/2 model, deriving the stochastic differential equations for the underlying asset price and the instantaneous variance process under the risk-neutral measure. It derives the conditional distribution properties of the integrated variance and the moment generating function representation of the underlying price, clarifying the economic implications of the model's core parameters to lay the theoretical foundation for the subsequent construction and empirical calibration of pricing methods.

Chapter Four discusses Option Pricing Methods and Algorithmic Implementation. Within the theoretical framework of the 3/2 model, this chapter systematically constructs a complete implementation system for the nine option pricing methods across three paradigms. It derives the core theoretical formulas for time-stepping conditional Monte Carlo simulations, exact and near-exact terminal simulation methods, and Fourier time-domain pricing methods. It also provides the numerical implementation flow for each method and conducts theoretical analyses on their error characteristics, convergence properties, and applicable scenarios, establishing a unified methodological foundation for subsequent numerical experiments.

Chapter Five focuses on Numerical Experiment Setup and Comparative Analysis of Results. This chapter establishes a unified numerical experiment environment and evaluation metric system, designing seven sets of test cases covering both conventional and extreme parameter environments. It systematically completes numerical testing for the nine pricing methods, quantitatively comparing their numerical stability, pricing accuracy, and computational efficiency, ultimately forming a hierarchical system of method performance and scenario-based selection recommendations.

Chapter Six presents the Empirical Calibration and Pricing Testing of the 3/2 Stochastic Volatility Model in Extreme Cryptocurrency Markets. Taking the extreme Bitcoin market flash crash of October 2025 as the research sample, this chapter constructs a complete empirical framework of ``long-term fixed-parameter time-series estimation - time-varying parameter cross-sectional calibration - pricing method extreme-scenario re-evaluation.'' It identifies the time-varying patterns of the model's core parameters throughout the crash event and the model's fitting effectiveness, validates the model's ability to fit the evolution of the volatility smile during extreme conditions, and compares these findings horizontally with the academic benchmark results from Chapter Five, defining the applicability boundaries of different pricing methods in extreme market environments.

Chapter Seven covers Conclusions and Outlook. This chapter systematically summarizes the core research conclusions of the entire thesis, objectively analyzes the limitations of this study, and provides an outlook on future research directions from the dimensions of model extension, method optimization, empirical scope expansion, and risk management applications.

Motivated by the operational failures and severe latencies of standard pricing algorithms under boundary conditions, this thesis seeks to overhaul the valuation framework for the $3/2$ stochastic volatility model.

Our research trajectory is anchored in two primary phases. First, we reconstruct the five prevailing numerical schemes, benchmarking their performance against vanilla options across a diverse grid of market regimes. By dissecting the resulting simulation data, the study maps out the exact compromise each algorithm makes between computational throughput and pricing fidelity, providing a transparent view of their operational boundaries.

The core contribution of this work, however, is a methodological breakthrough designed specifically to bypass the severe limitations these traditional models face with at-the-money (ATM) and short-tenor options. To resolve these bottlenecks, we introduce a highly robust pricing mechanism. This advanced approach consistently delivers rapid, high-precision valuations, retaining its structural integrity regardless of whether the market environment is tranquil or under extreme stress.